\hypertarget{appendix-932cf7182ff1}{}
\section{Phase 8 --- Linear Algebra: Active Recall}\label{appendix-932cf7182ff1-phase-8-linear-algebra-active-recall}

\textbf{Fast spine:} Measurements constrain → rank separates knowledge from freedom → factorizations expose structure → eigenvectors describe square-map directions → SVD describes every map → stable APIs report evidence.

Read the full explanations in \href{../01-Foundations-I/main.pdf\#nameddest=01-foundations-i-ch019-the-graphics-problem-hiding-inside-a-matrix-equation}{From Constraints to Spectra (01-Foundations-I, 01-foundations-i-ch019-the-graphics-problem-hiding-inside-a-matrix-equation)}. These figure links jump to their source chapters.

\subsection{The story at a glance}\label{appendix-932cf7182ff1-the-story-at-a-glance}

\pandocbounded{\includegraphics[keepaspectratio,alt={Diagram}]{../current/Assets/Print/b02539ebd0fde2394adee8c7dd9e84ffff1296acd8fd940824420640f8eda1a2.pdf}}

\subsection{Retrieve the mechanisms}\label{appendix-932cf7182ff1-retrieve-the-mechanisms}

{\def\LTcaptype{none} % do not increment counter
\begin{longtable}[]{@{}
  >{\raggedright\arraybackslash}p{(\linewidth - 4\tabcolsep) * \real{0.3333}}
  >{\raggedright\arraybackslash}p{(\linewidth - 4\tabcolsep) * \real{0.3333}}
  >{\raggedright\arraybackslash}p{(\linewidth - 4\tabcolsep) * \real{0.3333}}@{}}
\toprule\noalign{}
\begin{minipage}[b]{\linewidth}\raggedright
Cue
\end{minipage} & \begin{minipage}[b]{\linewidth}\raggedright
Short answer
\end{minipage} & \begin{minipage}[b]{\linewidth}\raggedright
Read / view
\end{minipage} \\
\midrule\noalign{}
\endhead
\bottomrule\noalign{}
\endlastfoot
Why may \(A\mathbf x=\mathbf b\) have no answer? & The target is outside \(\operatorname{Col}(A)\). & \href{../01-Foundations-I/main.pdf\#nameddest=01-foundations-i-ch019-a-linear-equation-is-a-measurement-and-several-equations-are-geometry}{Constraint geometry (01-Foundations-I, 01-foundations-i-ch019-a-linear-equation-is-a-measurement-and-several-equations-are-geometry)} · \href{../01-Foundations-I/main.pdf\#nameddest=01-foundations-i-ch019-figure-l2---consistency-freedom-and-the-four-fundamental-spaces}{L2 (01-Foundations-I, 01-foundations-i-ch019-figure-l2-\/-\/-consistency-freedom-and-the-four-fundamental-spaces)} \\
What does a null vector mean? & An input change no measurement can see; add it to any solution. & \href{../01-Foundations-I/main.pdf\#nameddest=01-foundations-i-ch019-rank-separates-what-the-measurements-determine-from-what-they-leave-free}{Four spaces (01-Foundations-I, 01-foundations-i-ch019-rank-separates-what-the-measurements-determine-from-what-they-leave-free)} · \href{../01-Foundations-I/main.pdf\#nameddest=01-foundations-i-ch019-figure-l2---consistency-freedom-and-the-four-fundamental-spaces}{L2 (01-Foundations-I, 01-foundations-i-ch019-figure-l2-\/-\/-consistency-freedom-and-the-four-fundamental-spaces)} \\
What does elimination preserve? & The solution set, only if row operations also transform \(\mathbf b\). & \href{../01-Foundations-I/main.pdf\#nameddest=01-foundations-i-ch019-elimination-turns-the-geometry-into-an-auditable-calculation}{Elimination (01-Foundations-I, 01-foundations-i-ch019-elimination-turns-the-geometry-into-an-auditable-calculation)} · \href{../01-Foundations-I/main.pdf\#nameddest=01-foundations-i-ch019-figure-l3---elimination-stores-work-for-later-solves}{L3 (01-Foundations-I, 01-foundations-i-ch019-figure-l3-\/-\/-elimination-stores-work-for-later-solves)} \\
Why factor once? & Reuse \(L,U\) or \(Q,R\) for additional right-hand sides. & \href{../01-Foundations-I/main.pdf\#nameddest=01-foundations-i-ch019-figure-l3---elimination-stores-work-for-later-solves}{L3 (01-Foundations-I, 01-foundations-i-ch019-figure-l3-\/-\/-elimination-stores-work-for-later-solves)} \\
Which factorization for what structure? & LU general square; Cholesky SPD; QR least squares; SVD rank insight; iteration for huge sparse systems when suitable. & \href{../01-Foundations-I/main.pdf\#nameddest=01-foundations-i-ch020-structure-makes-a-decomposition-a-better-tool-than-generic-elimination}{Factorization guide (01-Foundations-I, 01-foundations-i-ch020-structure-makes-a-decomposition-a-better-tool-than-generic-elimination)} · \href{../01-Foundations-I/main.pdf\#nameddest=01-foundations-i-ch020-figure-l4---matrix-structure-determines-the-useful-factorization}{L4 (01-Foundations-I, 01-foundations-i-ch020-figure-l4-\/-\/-matrix-structure-determines-the-useful-factorization)} \\
What does a small residual prove? & Near satisfaction of the equations; sensitivity can still make the solution inaccurate. & \href{../01-Foundations-I/main.pdf\#nameddest=01-foundations-i-ch020-eigenvectors-reveal-the-directions-a-square-map-does-not-turn}{Eigenpair sensitivity (01-Foundations-I, 01-foundations-i-ch020-eigenvectors-reveal-the-directions-a-square-map-does-not-turn)} · \href{../01-Foundations-I/main.pdf\#nameddest=01-foundations-i-ch019-matrix-norms-give-scale-to-residuals-and-error-claims}{Norms and error (01-Foundations-I, 01-foundations-i-ch019-matrix-norms-give-scale-to-residuals-and-error-claims)} \\
Why is a symmetric eigenbasis special? & It is orthonormal, preserving 2-norm while diagonalizing the map. & \href{../01-Foundations-I/main.pdf\#nameddest=01-foundations-i-ch020-figure-l5---eigenanalysis-is-cleanest-when-the-basis-is-orthonormal}{L5 (01-Foundations-I, 01-foundations-i-ch020-figure-l5-\/-\/-eigenanalysis-is-cleanest-when-the-basis-is-orthonormal)} \\
What is the SVD action? & Input orthogonal coordinates → independent stretches → output orthogonal coordinates. & \hyperref[02-foundations-ii-ch024-singular-value-decomposition-covers-every-rectangular-map]{SVD geometry} · \hyperref[02-foundations-ii-ch024-figure-l6---every-linear-map-is-orthogonal-change-stretch-orthogonal-change]{L6} \\
What does the pseudoinverse choose? & Exact \(A^+\): least-squares fit, then minimum norm. Truncated \(A_\tau^+\): solve the retained-rank problem; original residual can increase. & \hyperref[02-foundations-ii-ch024-pseudoinverse-makes-least-squares-and-minimum-norm-one-operation]{Pseudoinverse} · \hyperref[02-foundations-ii-ch024-figure-l7---the-singular-spectrum-makes-inversion-and-compression-explicit]{L7} \\
How does PCA connect to SVD? & Center rows, factor \(X=U\Sigma V^\mathsf T\); \(V\) gives principal directions and \(U\Sigma\) scores. & \hyperref[02-foundations-ii-ch024-pca-turns-a-cloud-of-points-into-measured-variation]{PCA} \\
What does polar decomposition separate? & Orthogonal motion from symmetric stretch; reflection and singular ambiguity need attention. & \hyperref[02-foundations-ii-ch024-polar-decomposition-separates-rotation-like-motion-from-stretch]{Polar decomposition} \\
Which SVD checks are independent? & Reconstruction, \(U\) orthogonality, \(V\) orthogonality, ordering, dimensions, convergence. & \hyperref[02-foundations-ii-ch038-figure-l8---svd-validation-needs-independent-checks]{L8} \\
When can diagonalization explain a matrix function? & When its eigenbasis is usable; otherwise use a stable function-specific algorithm. & \hyperref[02-foundations-ii-ch026-matrix-functions-turn-repeated-linear-actions-into-one-operator]{Matrix functions} · \hyperref[02-foundations-ii-ch026-figure-l9---a-matrix-function-follows-the-same-action-after-a-change-of-basis]{L9} \\
Direct or iterative solve? & Decide from size, sparsity, right-hand-side reuse, convergence, memory, and error needs. & \hyperref[02-foundations-ii-ch038-figure-l10---direct-and-iterative-methods-occupy-different-workloads]{L10} \\
\end{longtable}
}

\subsection{Thirty-second reconstruction}\label{appendix-932cf7182ff1-thirty-second-reconstruction}

\begin{enumerate}
\def\labelenumi{\arabic{enumi}.}
\tightlist
\item
  Sketch one measurement as a hyperplane; explain consistency using column space.
\item
  Explain why rank and nullity count determined and invisible directions.
\item
  Derive one pivot step and show the same row operation on the right-hand side.
\item
  Select LU, Cholesky, QR, SVD, or an iterative method from matrix structure.
\item
  Explain why normal equations square the condition number.
\item
  Trace \(A=U\Sigma V^\mathsf T\) on one vector.
\item
  State what thresholding changes in a pseudoinverse or truncated SVD.
\item
  Name separate SVD correctness checks and what a residual cannot prove.
\end{enumerate}

\href{../01-Foundations-I/main.pdf\#nameddest=01-foundations-i-ch019-the-graphics-problem-hiding-inside-a-matrix-equation}{Full teaching volume (01-Foundations-I, 01-foundations-i-ch019-the-graphics-problem-hiding-inside-a-matrix-equation)} · \href{../01-Foundations-I/main.pdf\#nameddest=appendix-fd090d982a15-phase-8---parts-701-800}{100-part title ledger (01-Foundations-I, appendix-fd090d982a15-phase-8-\/-\/-parts-701-800)} · \hyperref[appendix-4ad47bdd269a-phase-8---from-constraints-to-spectra]{Dependency spine}

\subsection{Fast retrieval: restore the details that are easy to blur}\label{appendix-932cf7182ff1-fast-retrieval-restore-the-details-that-are-easy-to-blur}

\hyperref[02-foundations-ii-ch024-figure-l6a---watch-the-same-directions-become-an-ellipse]{Figure L6a: exact SVD geometry} · \href{../current/Assets/Print/3491f9bdb1d7f300619c0d11a33ded7f3b3adfdc803929f353124f6903c784b4.svg}{Editable SVG}

\begin{figure}
\centering
\pandocbounded{\includegraphics[keepaspectratio,alt={Visual Recall/phase8\_svd\_geometry.png}]{../current/Assets/Print/b681e335dbaa5c1c27f9d47821582f982e787aff8b237276782bfc38d8631fb1.png}}
\caption{Visual Recall/phase8\_svd\_geometry.png}
\end{figure}

{\def\LTcaptype{none} % do not increment counter
\begin{longtable}[]{@{}
  >{\raggedright\arraybackslash}p{(\linewidth - 4\tabcolsep) * \real{0.3333}}
  >{\raggedright\arraybackslash}p{(\linewidth - 4\tabcolsep) * \real{0.3333}}
  >{\raggedright\arraybackslash}p{(\linewidth - 4\tabcolsep) * \real{0.3333}}@{}}
\toprule\noalign{}
\begin{minipage}[b]{\linewidth}\raggedright
Cue
\end{minipage} & \begin{minipage}[b]{\linewidth}\raggedright
Shortest reliable answer
\end{minipage} & \begin{minipage}[b]{\linewidth}\raggedright
Deep read / figure
\end{minipage} \\
\midrule\noalign{}
\endhead
\bottomrule\noalign{}
\endlastfoot
When does Gram--Schmidt normalize a new column? & Only if its residual norm is safely nonzero; rank deficiency needs rank handling or completion. & \href{../01-Foundations-I/main.pdf\#nameddest=01-foundations-i-ch020-orthogonal-coordinates-preserve-the-shape-of-the-problem}{Gram--Schmidt (01-Foundations-I, 01-foundations-i-ch020-orthogonal-coordinates-preserve-the-shape-of-the-problem)} \\
What does a Householder sign choice prevent? & Choose \(\alpha=-\operatorname{copysign}(\|x\|,x_1)\) to avoid cancellation; compute the norm safely. & \href{../01-Foundations-I/main.pdf\#nameddest=01-foundations-i-ch020-reflections-and-plane-rotations-build-qr-without-losing-lengths}{Householder reflectors (01-Foundations-I, 01-foundations-i-ch020-reflections-and-plane-rotations-build-qr-without-losing-lengths)} \\
When is Givens rotation undefined? & If \(a=b=0\), skip it; otherwise use \texttt{hypot(a,b)} for the scale. & \href{../01-Foundations-I/main.pdf\#nameddest=01-foundations-i-ch020-reflections-and-plane-rotations-build-qr-without-losing-lengths}{Givens rotations (01-Foundations-I, 01-foundations-i-ch020-reflections-and-plane-rotations-build-qr-without-losing-lengths)} \\
What residual remains after QR least squares? & The irreducible perpendicular part \(b_\perp=(I-QQ^T)b\); minimum residual norm is \(\|b_\perp\|\). & \href{../01-Foundations-I/main.pdf\#nameddest=01-foundations-i-ch020-qr-turns-least-squares-into-a-triangular-solve-plus-an-irreducible-residual}{QR residual derivation (01-Foundations-I, 01-foundations-i-ch020-qr-turns-least-squares-into-a-triangular-solve-plus-an-irreducible-residual)} \\
What makes an eigenvector estimate sensitive? & A small residual is backward evidence; a small eigengap can still rotate the vector or subspace. & \href{../01-Foundations-I/main.pdf\#nameddest=01-foundations-i-ch020-residuals-distinguish-backward-accuracy-from-forward-certainty}{Residual and eigengap (01-Foundations-I, 01-foundations-i-ch020-residuals-distinguish-backward-accuracy-from-forward-certainty)} · \href{../01-Foundations-I/main.pdf\#nameddest=01-foundations-i-ch020-figure-l5---eigenanalysis-is-cleanest-when-the-basis-is-orthonormal}{L5 (01-Foundations-I, 01-foundations-i-ch020-figure-l5-\/-\/-eigenanalysis-is-cleanest-when-the-basis-is-orthonormal)} \\
What does a real Schur 2×2 block mean? & A real representation of a complex-conjugate eigenvalue pair; do not read only its diagonal. & \href{../01-Foundations-I/main.pdf\#nameddest=01-foundations-i-ch020-qr-iteration-and-jacobi-rotations-preserve-the-spectrum-while-simplifying-shape}{QR and Schur form (01-Foundations-I, 01-foundations-i-ch020-qr-iteration-and-jacobi-rotations-preserve-the-spectrum-while-simplifying-shape)} \\
Where is the plane-fit normal in an \(N\times3\) data SVD? & In the right singular vectors \(V\); take the smallest-variance direction. & \hyperref[02-foundations-ii-ch024-plane-fits-eigengaps-and-interpretation-of-uncertain-components]{Plane fit and eigengap} \\
What must a wide SVD return? & With \(r=\min(m,n)\), thin \(U:m\times r\), \(\sigma:r\), \(V^T:r\times n\); factor \(A^T\) if using a tall backend. & \hyperref[02-foundations-ii-ch038-householder-reductions-turn-a-dense-svd-into-a-bidiagonal-problem]{Wide SVD shapes} · \hyperref[02-foundations-ii-ch024-figure-l6---every-linear-map-is-orthogonal-change-stretch-orthogonal-change]{L6} \\
What are zero singular vectors for? & Right zero directions span \(\operatorname{Null}(A)\); full \(U\) zero/completion directions span \(\operatorname{Null}(A^T)\). & \hyperref[02-foundations-ii-ch038-zero-singular-directions-require-deliberate-orthonormal-completion]{Zero-singular completion} \\
What is the exact \(A^+\) rule? & Reciprocate every positive \(\sigma_i\); it minimizes original residual, then chooses minimum norm. & \hyperref[02-foundations-ii-ch024-the-moorepenrose-inverse-and-its-four-identities]{Exact pseudoinverse} · \hyperref[02-foundations-ii-ch024-figure-l7---the-singular-spectrum-makes-inversion-and-compression-explicit]{L7} \\
What does \(A_\tau^+\) change? & It discards positive values at or below \(\tau\); the original least-squares residual can increase. & \hyperref[02-foundations-ii-ch024-thresholding-and-regularization-are-different-choices]{Thresholding versus regularization} · \hyperref[02-foundations-ii-ch024-figure-l7---the-singular-spectrum-makes-inversion-and-compression-explicit]{L7} \\
What errors does rank-\(k\) SVD leave? & Spectral error \(\sigma_{k+1}\); Frobenius squared error \(\sum_{i>k}\sigma_i^2\). & \hyperref[02-foundations-ii-ch024-low-rank-approximation-and-why-the-largest-directions-win]{Eckart--Young errors} \\
How do PCA scores relate to SVD? & For centered row samples \(X=U\Sigma V^T\), scores are \(XV=U\Sigma\); covariance eigenvalues are \(\sigma_i^2/(N-1)\). & \hyperref[02-foundations-ii-ch024-centered-measurements-reveal-principal-directions]{PCA scores} \\
What does polar decomposition call a reflection? & \(UV^T\) is orthogonal; if its determinant is \(-1\), it is improper, and nearest \(SO(n)\) flips a least-singular direction. & \hyperref[02-foundations-ii-ch024-two-sided-rotations-nearest-proper-rotation-and-nonunique-directions]{Nearest proper rotation} \\
Which Jacobi angle convention is in use? & For \(J=[[c,s],[-s,c]]\), use \(p'=cp-sq\), \(q'=sp+cq\), \(\tau=(\beta-\alpha)/(2\gamma)\); skip \(\gamma=0\). & \hyperref[02-foundations-ii-ch038-jacobi-svd-removes-column-coupling-pair-by-pair]{Jacobi convention} \\
What bounds a strided matrix view? & Check \(c\le s\), checked last offset \((r-1)s+(c-1)<e\), and null pointer only with zero extent. & \hyperref[02-foundations-ii-ch038-matrix-views-need-checked-address-arithmetic-as-well-as-a-stride]{View bounds} · \hyperref[02-foundations-ii-ch038-workspace-allocation-is-part-of-the-solver-contract]{Workspace bounds} \\
When is the principal matrix log defined? & When the spectrum avoids the closed negative real axis, including zero; other real logs may still exist. & \hyperref[02-foundations-ii-ch026-logarithms-and-square-roots-need-domains-and-branches]{Log domain and branches} · \hyperref[02-foundations-ii-ch026-figure-l9---a-matrix-function-follows-the-same-action-after-a-change-of-basis]{L9} \\
What must CG preconditioning preserve? & Use SPD PCG/symmetric preconditioning; vanilla CG on arbitrary \(M^{-1}A\) may lose symmetry. Stop before zero-residual divisions. & \hyperref[02-foundations-ii-ch038-direct-and-iterative-solvers-suit-different-workloads]{CG preconditioning and true residual} · \hyperref[02-foundations-ii-ch038-figure-l10---direct-and-iterative-methods-occupy-different-workloads]{L10} \\
What makes an SVD API result inspectable? & Shapes/mode, rank threshold, workspace, convergence status, and independent reconstruction/orthogonality checks. & \hyperref[02-foundations-ii-ch038-an-svd-api-begins-with-shapes-and-guarantees]{SVD API contract} · \hyperref[02-foundations-ii-ch038-figure-l8---svd-validation-needs-independent-checks]{L8} \\
\end{longtable}
}

\begin{quote}
Read the linked teaching sections for the derivations and caveats; this table is a retrieval prompt, not a replacement explanation.
\end{quote}

\subsection{Added retrieval prompt}\label{appendix-932cf7182ff1-added-retrieval-prompt}

Whiten by solving Lz=x-mu; then Mahalanobis squared distance is z transpose z. A singular covariance needs lower-dimensional support, not full-area density. \hyperref[02-foundations-ii-ch024-original-bridge-whiten-a-covariance-ellipse]{Rebuild the example}.

\begin{figure}
\centering
\pandocbounded{\includegraphics[keepaspectratio,alt={Visual Recall/Math Foundations/M08-1\_cholesky-whitening.png}]{../current/Assets/Print/e4e424d7b45ba80ab1141d375568bed99279bbe0965c9654872418946ee0692d.png}}
\caption{Visual Recall/Math Foundations/M08-1\_cholesky-whitening.png}
\end{figure}

\textbf{M08-1. A covariance ellipse and its whitened coordinates.} For Sigma=\texttt{{[}{[}4,1.2{]},{[}1.2,1{]}{]}}, L=\texttt{{[}{[}2,0{]},{[}0.6,0.8{]}{]}}, mu=(3,-1), z=(1,1), x=(5,0.4), Mahalanobis squared distance 2, Euclidean squared distance 5.96, and density exp(-1)/(3.2 pi), about 0.0366 for an SPD 2D Gaussian.
